\documentclass[11pt, oneside]{article} 
\usepackage{geometry}
\geometry{letterpaper} 
\usepackage{graphicx}
	
\usepackage{amssymb}
\usepackage{amsmath}
\usepackage{parskip}
\usepackage{color}
\usepackage{hyperref}

\graphicspath{{/Users/telliott/Dropbox/Github-math/figures/}}
% \begin{center} \includegraphics [scale=0.4] {gauss3.png} \end{center}

\title{Pizza theorem}
\date{}

\begin{document}
\maketitle
\Large

%[my-super-duper-separator]

Here is another problem that I saw in Acheson's book:  it is called the "pizza theorem".  

\begin{center} \includegraphics [scale=0.6] {Acheson_G111.png} \end{center}

Choose a point \emph{anywhere} in the pie, draw two chords at right angles and then fill in with two more chords at $45^{\circ}$.

Form the sum of the areas of alternate slices.  The total dark area and the total light area are always equal.

When one of the chords is a diameter of the circle, this result is obviously correct, by symmetry.  The goal will be to prove that it is true even when we rotate the perpendicular chords about the point.

\subsection*{Extraordinary property}

\label{sec:extraordinary_property}

As a preliminary result, let's do one more of his problems:  it comes from a book by Malton and is described as an "extraordinary property of the circle".

\begin{center} \includegraphics [scale=0.6] {Acheson_G110.png} \end{center}

$\bullet$  Let two chords of a circle meet at right angles.  Then the squares of the four components add up to a constant, and that constant is equal to $(2R)^2$, twice the radius of the circle, squared.

The key to the proof is to draw both diameters of the circle that terminate on each end of one of the chords, let's say the one composed of $a + b$.

\begin{center} \includegraphics [scale=0.4] {perp_chords3.png} \end{center}

If the four points at the ends of the two diameters are joined to form a rectangle (only the horizontal components are shown in the figure above), then it is clear by mirror-image symmetry that the short extension at the top is also equal to $c$, so the height of the rectangle is $d - c$.

We will prove this more formally, using the following preliminary results.

The first is that when any two diameters of a circle are drawn and the points on the circle joined together, the result is a rectangle.  We proved this \hyperref[sec:rectangle_in_a_circle]{\textbf{here}}.

\begin{center} \includegraphics [scale=0.4] {perp_chords4.png} \end{center}

Using this result, find the diameters terminating at $S$ and $R$ (not shown), and then draw the rectangle $PQRS$.

The original vertical chord is shown, which is divided into two segments $c$ and $d$ ($d$ is not labeled).  The vertical chord is given as perpendicular to $SR$.  

Draw radius $OT$ perpendicular to the vertical chord.  By our work \hyperref[sec:perpendicular_bisector_of_a_chord]{\textbf{here}}, $OT$ bisects both the side of the rectangle $PS$ and the vertical chord.

\begin{center} \includegraphics [scale=0.4] {perp_chords5.png} \end{center}

Bisection of $PS$, together with all the right angles, means that the two small rectangles with sides labeled $a$ and $e$ are congruent.  And since the chord was bisected, the line segment at the top of the figure also has length $c$, since $c + e$ is the same length above and below the bisector.

\begin{center} \includegraphics [scale=0.4] {perp_chords3.png} \end{center}

$\square$

Now for the main theorem.

\emph{Proof.}

This is trivial, now that we have $d - c$.  By the Pythagorean theorem
\[ (a + b)^2 + (d - c)^2 = (2R)^2 \]
\[ = a^2 + b^2 + d^2 + c^2 + 2ab - 2cd \]

But $ab = cd$ so
\[ 4R^2 = a^2 + b^2 + c^2 + d^2 \]

$\square$

Acheson says that this theorem leads to a proof of the "pizza theorem".  

\begin{center} \includegraphics [scale=0.6] {Acheson_G111.png} \end{center}

\subsection*{plausability proof}

To begin with, it is clear that if the chosen point is at the center, then all the segments are congruent.  Furthermore, if the point lies on a diagonal, there will be four different types of segments, each type with one light and one dark example, (red dots have been added to show to which group the segments belong in the figure below).

\begin{center} \includegraphics [scale=0.4] {pizza3.png} \end{center}

Here is a general point, not on a diagonal.

\begin{center} \includegraphics [scale=0.4] {pizza4.png} \end{center}

Any such general point in the circle can be reached by moving off the diagonal in a simple way.  

One is a horizontal translation.

\begin{center} \includegraphics [scale=0.4] {pizza5.png} \end{center}

Another way to move one of the boundaries onto or off of a diagonal is by rotation.

\begin{center} \includegraphics [scale=0.4] {pizza6.png} \end{center}

We will show that rotation doesn't change the total area of slices marked with red dots.  This is the operation that is connected to the special property of the circle which we gave above.

\subsection*{area of a sector}

For a small angle $\Delta \theta$, the area of a sector swept out by that angle is 
\[ A = \frac{1}{2} \Delta \theta \ r \cdot r = \Delta \theta \ \frac{r^2}{2} \]

If we look at the sectors marked with red dots, a counter-clockwise rotation adds to the area on the left, viewed from the point and facing out.  It will be seen that the relevant radii are all at right angles to one another.

These values of $r$ are exactly those which were involved in the previous proof:  one is $\Delta \theta a^2/2$, another has $b^2$, and so on (using the notation from the special property example).

Therefore, the total increase in area is
\[ \Delta A = \Delta \theta \ \frac{a^2 + b^2 + c^2 + d^2}{2} \]
\[ = \Delta \theta \ \frac{4R^2}{2} = \Delta \theta \ 2R^2 \]

The increase in area for a small rotation is a constant.  

The unshaded areas, although they have different radii, add exactly the same area, and for the same reason.

The area gained by white is the same as the area gained by dark, but it is also \emph{precisely} the same area, (not just quantitativlely, but factually) as that lost by dark.

For us at present, the result is only plausible and not rigorous, because we have assumed that $r$ doesn't change with $\Delta \theta$.  It does change, slowly, and importantly, these changes balance each other.  For two shaded arcs the radius is increasing, and for the other two decreasing, so that the net change is zero.

We need calculus to do more than this.

\url{https://www.math.uni-bielefeld.de/~sillke/PUZZLES/pizza-theorem}

In the language of calculus we would have that (twice) the area added by a rotation is

\[ \sum_i r_i^2 (\phi_i) \cdot \Delta \phi = \sum_i  \int_{\theta_1}^{\theta_2} r_i^2 (\phi_i) \ d \phi \]

But the sum of the integrals is the integral of the sum.  

\[ = \int_{\theta_1}^{\theta_2}  \sum_i  r_i^2 (\phi_i) \ d \phi \]

And when we add them together, the sum of those factors $r^2$ is a constant, by the extraordinary property of the circle.  So in the end we obtain $(2R)^2$ times the total of the four angular rotations.  The changes in the radius simply drop out from the calculation.

\subsection*{references}

I found a number of images for a "proof without words" online and I also found a reference

to some related proofs, as well as a discussion on Stack Exchange:

\url{https://math.stackexchange.com/questions/865818/can-anyone-explain-pizza-theorem}

The reference says

\begin{quote}Sliding these arcs and
chords together, we see that the chords form a right triangle with the
diameter of the circle as the hypotenuse.\end{quote}

The various pictures of the proof without words do not explain how the shapes were arrived at.  However, they do show that the complementary light and dark areas each add up to one-half of the pizza.

\begin{center} \includegraphics [scale=0.4] {pizza2.png} \end{center}

\end{document}